\lysh{15271_SOLUTION}{1}
\textbf{ΛΥΣΗ}

\(\alpha\)) Η ευθεία που διέρχεται από τα σημεία Α, Β έχει συντελεστή διεύθυνσης $\lambda = \frac{6-2}{1+3} = 1$.

\(\beta\)) Η ευθεία έχει συντελεστή διεύθυνσης $\lambda = 1$ και διέρχεται από το σημείο Α, οπότε έχει εξίσωση $y-2=1(x+3)$ δηλαδή είναι η $y=x+5$.

\(\gamma\)) Το σημείο $\Gamma$ είναι πάνω στην ευθεία ΑΒ μόνο όταν οι συντεταγμένες του επαληθεύουν την εξίσωση της ευθείας ΑΒ.
Με $x=-13$ στον τύπο της ευθείας ΑΒ έχουμε:
\[ y=-13+5=-8 \neq y_{\Gamma}=-7 \]
Άρα το σημείο $\Gamma$ δεν είναι πάνω στην ΑΒ.
\end{lysh}